\documentclass[10pt,a4paper]{amsart}
\usepackage[margin=19mm]{geometry}
\usepackage{amsmath,amssymb,mathtools}
\usepackage[scr=rsfs,cal=euler]{mathalfa}
\usepackage{xfrac}
\usepackage[unicode,hidelinks,bookmarks=false]{hyperref}
\newcommand{\ie}{i.e.\ }
\newcommand{\cf}{cf.\ }
\newcommand{\resp}{respectively\ }
\newcommand{\Subsection}{\subsection}
\providecommand{\nobreakdash}{\nobreak\hbox{-}}
\setlength{\emergencystretch}{2em}
\multlinegap0pt
\usepackage{bm}
\usepackage{bbm}
\usepackage{dsfont}
\usepackage{xspace}
\usepackage{cancel}
\usepackage{mathtools}
\usepackage{amsthm}
\makeatletter
\@ifundefined{lemma}{\newtheorem{lemma}{Lemma}}{}
\@ifundefined{theorem}{\newtheorem{theorem}{Theorem}}{}
\makeatother
\title[Critical perturbation sizes for preserving decay rates in so]{Critical perturbation sizes for preserving decay rates in solutions of perturbed nonlinear differential equations}
\author{John A. D. Appleby, Subham Pal}
\date{Source: arXiv:2507.10151v3 (CC BY 4.0)}
\begin{document}
\maketitle

\noindent\emph{Source and licence.} Critical perturbation sizes for preserving decay rates in solutions of perturbed nonlinear differential equations. Version arXiv:2507.10151v3 (CC BY 4.0), distributed under the Creative Commons Attribution 4.0 International licence, as recorded in the arXiv licence metadata for this exact version and shown on the abstract page https://arxiv.org/abs/2507.10151v3. Full rights notice: licenses/arxiv-2507.10151-CC-BY-4.0.txt.\par\smallskip

\noindent\textit{Editorial scope.} Counted unit: the Theorem whose source label is \texttt{thm.Xwithouttailmart} in the article (source0.tex lines 2675--2696, source label \texttt{thm.Xwithouttailmart}), delivered together with the article's own complete original proof of it (source0.tex lines 2697--2838, one \texttt{proof} environment of 142 lines). No mathematics is written, repaired or re-derived: statement and proof are the article's own text line for line. Required context retained uncounted: eq.odepert (lines 90--92); eq.fglobalstable (lines 104--106); def.F (lines 125--127); eq.fasypres (lines 157--159); eq.fasypresmu (lines 200--202); eq.g2 (lines 1082--1084); eq:Gamma (lines 1086--1088); eq.GammaOFinv (lines 1138--1140); thm.bigOconverse (lines 2146--2161); eq.sde7 (lines 2405--2407); thm.Xwithtailmart (lines 2431--2441); lemma.tailmartasy (lines 2540--2553); d (lines 2541--2543). Every definition, display and label of those lines is retained unchanged. Omitted from this package and not redistributed: 1--89, 93--103, 107--124, 128--156, 160--199, 203--1081, 1085--1085, 1089--1137, 1141--2145, 2162--2404, 2408--2430, 2442--2539, 2544--2674, 2839--3076. Nothing inside a retained excerpt is omitted or altered: every retained body is delivered exactly as the article's lines are written, so no omission receipt of any kind accompanies this package. Citations: the retained excerpts carry no citation command at all, so no bibliography entry of the article is transcribed and no reference is redistributed. Reference adaptation: none was needed and none was made; every reference inside the delivered unit resolves inside it, so no unresolved reference or citation is produced. Printed slips kept literal: no printed word, symbol, hypothesis or number of the counted statement or of its proof was altered. Layout and class: the article's own class is \texttt{amsart} with packages including amssymb, latexsym, bbm, amsmath, amsfonts, amsthm, eucal, rotating, multirow; the wrapper is the lane's pinned \texttt{amsart} editorial wrapper at 10pt on A4, which restates the theorem-like environments the retained text needs and adds the article's own macro definitions that the retained text needs (none), with extra wrapper packages (bm, bbm, dsfont, xspace, cancel, mathtools). The delivered numbering is the unit's own. Assets: no retained span contains a figure environment, a graphics-inclusion command, a PDF-page inclusion, a table or a TikZ picture; no image file is copied into this package, the package declares no asset dependency and no image file is redistributed with it. Licence: version arXiv:2507.10151v3 of this article is distributed under the Creative Commons Attribution 4.0 International licence, as recorded in the arXiv licence metadata for this exact version and shown on the abstract page https://arxiv.org/abs/2507.10151v3. A full rights notice is delivered at licenses/arxiv-2507.10151-CC-BY-4.0.txt. Characters: every retained excerpt is pure ASCII, so the delivered engine, XeLaTeX with its default Latin Modern text font, reports no missing character.
	\begin{equation} \label{eq.odepert}
		x'(t)=-f(x(t))+g(t), \quad t>0; \quad x(0)=\xi,
	\end{equation}

	\begin{equation} \label{eq.fglobalstable}
		xf(x)>0 \quad\text{for $x\neq 0$,} \quad f(0)=0.
	\end{equation}

	\begin{equation} \label{def.F}
		F(x)=\int_x^1 \frac{1}{f(u)}\,du, \quad x>0,
	\end{equation}

	\begin{equation} \label{eq.fasypres}
		\liminf_{x\to 0^+} \frac{f((1-\epsilon)x)}{f(x)}= \underline{\Phi}_f(\epsilon)
	\end{equation}

	\begin{equation} \label{eq.fasypresmu}
		\limsup_{x\to 0^+} \frac{f(\mu x)}{f(x)}=:\bar{\Phi}_f(\mu)<\mu, \quad \text{ for all $\mu<1$}.
	\end{equation}

	\begin{equation} \label{eq.g2}
		\lim_{t\to\infty} \int_0^t g(s)\,ds \text{ exists and is finite}
	\end{equation}

	\begin{equation} \label{eq:Gamma}
		\Gamma(t):=-\int_t^\infty g(s)\,ds, \quad t\geq 0
	\end{equation}

	\begin{equation} \label{eq.GammaOFinv}
		\Gamma(t)=O(F^{-1}(t)), \quad t\to\infty,
	\end{equation}

	\begin{theorem} \label{thm.bigOconverse}
		Suppose that $f\in C(\mathbb{R};\mathbb{R})$ is increasing, odd and obeys \eqref{eq.fglobalstable}, \eqref{eq.fasypres}, and \eqref{eq.fasypresmu}. Suppose $F$ is defined by \eqref{def.F}.  Let $g$ be a continuous function. If the continuous solution of \eqref{eq.odepert}  obeys 
		\begin{align*}
			x(t)=O(F^{-1}(t)), \quad t\to\infty,
		\end{align*}	
		then $g$ is such that \eqref{eq.g2} holds and $\Gamma$ be defined by \eqref{eq:Gamma} obeys \eqref{eq.GammaOFinv}.  
		%and define
		%\[
		%F(x) := \int_{x}^{1} \frac{1}{f(u)} \, du, 
		%\qquad x>0 .
		%\]
		%Then $F$ is decreasing and invertible, with derivative
		%\[
		%F'(u) = -\frac{1}{f(u)} .
		%\]
	\end{theorem}	

	\begin{equation}  \label{eq.sde7}
		dX(t) = -f\big(X(t)\big)\,dt + \sigma(t)\,dB(t), \quad t\geq 0.
	\end{equation}

	\begin{theorem} \label{thm.Xwithtailmart}
		Suppose that $f\in C(\mathbb{R};\mathbb{R})$ is increasing, odd and obeys \eqref{eq.fglobalstable}, \eqref{eq.fasypres}, and \eqref{eq.fasypresmu}. Suppose $F$ is defined by \eqref{def.F}.  Let $\sigma$ be a continuous function, and let $X$ be the unique continuous adapted process satisfying \eqref{eq.sde7}.
		Then the following are equivalent:
		\begin{enumerate}
			\item[(A)] $\sigma \in L^2(0,\infty)$, \quad $\lim_{t \to \infty} \dfrac{\int_{t}^{\infty} \sigma(s)\, dB(s)}{F^{-1}(t)} = 0$, \; a.s.
			\item[(B)] There is a random variable $\Lambda$, which takes only the values $1,0,-1$, such that
			\[
			\lim_{t \to \infty} \frac{X(t)}{F^{-1}(t)} = \Lambda, \quad \text{a.s.}
			\]
		\end{enumerate}
	\end{theorem}

	\begin{lemma} \label{lemma.tailmartasy} Assume $\sigma$ is continuous, $\sigma \in L^2(0, \infty)$ and obeys
		\begin{equation} \label{d}
			\int_t^\infty \sigma^2(s) \, ds > 0, \quad \forall t \geq 0.
		\end{equation}
		Then,
		\begin{multline*}
			\limsup_{t \to \infty} \frac{\int_t^\infty \sigma(s) \, dB(s)}{\sqrt{2 \int_t^\infty \sigma^2(s) \, ds \log_2 \left( \frac{1}{\int_t^\infty \sigma^2(s) \, ds} \right)}} = 1\\
			= -\liminf_{t \to \infty} \frac{ \int_t^\infty \sigma(s) \, dB(s)}{\sqrt{2 \int_t^\infty \sigma^2(s) \, ds \log_2 \left( \frac{1}{\int_t^\infty \sigma^2(s) \, ds} \right)}}, \quad \text{a.s.}
		\end{multline*}
		where
		\[
		\log_2 x := \log(\log(x)).
		\]
	\end{lemma}

	\begin{theorem}  \label{thm.Xwithouttailmart}.
		Suppose that $f\in C(\mathbb{R};\mathbb{R})$ is increasing, odd and obeys \eqref{eq.fglobalstable}, \eqref{eq.fasypres}, and \eqref{eq.fasypresmu}. Suppose $F$ is defined by \eqref{def.F}. 
		Suppose $\sigma$ is continuous, $\sigma \in L^2(0,\infty)$ and obeys \eqref{d}, and that there exists $\mu \in [0,\infty]$ such that
		\[
		\mu := \lim_{t \to \infty} \frac{\int_t^\infty \sigma^2(s)\, ds }{F^{-1}(t)^2/\log\log(1/F^{-1}(t)^2)}\in [0,\infty]
		\]
		is well--defined. Then:
		\begin{itemize}
			\item[(i)] If $\mu = 0$, then there exists a $\mathcal{F}^B(\infty)$--measurable random variable $\lambda$ with $\mathbb{P}[\lambda\in \{-1,0,1\}]=1$ such that 
			\[
			\mathbb{P}\left[ \lim_{t \to \infty} \frac{X(t)}{F^{-1}(t)} =\lambda \right] = 1.
			\]
			\item[(ii)] If $\mu \in (0,\infty)$, then
			\[
			\mathbb{P}\left[ \lim_{t \to \infty} \frac{X(t)}{F^{-1}(t)} \in \{-1,0,1\} \right] = 0.
			\]
			\item[(iii)] If $\mu=+\infty$, then 
			\[
			\mathbb{P}\left[ \limsup_{t \to \infty} \frac{|X(t)|}{F^{-1}(t)} =+\infty \right] = 1.
			\]
		\end{itemize}
	\end{theorem}

	\begin{proof}
		We start by noticing that the asymptotic estimate on $\int_t^\infty \sigma^2(s)\,ds$ is equivalent to 
		\[
		\mu := \lim_{t \to \infty} \frac{\int_t^\infty \sigma^2(s)\, ds \; \log\log \left( \frac{1}{\int_t^\infty \sigma^2(s) \, ds} \right)}{F^{-1}(t)^2} \in [0,\infty], 
		\]
		by virtue of the discussion before the theorem. 
		
		We prove part (i) first, in which $\mu=0$. Define $\varphi(x)=x\log_2(1/x)$ for $0<x<e{-e}$. On this interval, $\varphi$ is increasing and also $\lim_{x\to 0^+} \varphi(x)=0$.
		
		By Lemma~\ref{lemma.tailmartasy}, we have 
		\[
		\limsup_{t \to \infty} \frac{|\Gamma(t)|}{\sqrt{2\varphi\!\left(\int_t^\infty \sigma^2(s)\, ds\right)}} = 1,\quad \text{a.s.}
		\]
		Since $\mu = 0$, we have
		\[
		\lim_{t \to \infty} \frac{\varphi\!\left(\int_t^\infty \sigma^2(s)\, ds\right)}{F^{-1}(t)^2} = 0.
		\]
		Hence,
		\[
		\limsup_{t \to \infty} \frac{|\Gamma(t)|}{F^{-1}(t)}
		= \limsup_{t \to \infty} 
		\left\{ \frac{|\Gamma(t)|}{\sqrt{2\varphi\!\left(\int_t^\infty \sigma^2(s)\, ds\right)}}
		\times
		\sqrt{\frac{2\varphi\!\left(\int_t^\infty \sigma^2(s)\, ds\right)}{F^{-1}(t)^2}}
		\right\}
		= 0, \quad \text{a.s.}
		\]
		Therefore, we have by Theorem \ref{thm.Xwithtailmart} that
		\[
		\lim_{t \to \infty} \frac{X(t)}{F^{-1}(t)} = \lambda \in \{-1,0,1\}, 
		\quad \text{a.s.},
		\]
		and the proof is complete.
		
		We now prove part (ii). Let
		\[
		A = \left\{ \omega : \lim_{t \to \infty} \frac{X(t,\omega)}{F^{-1}(t)} = \lambda(\omega) \in \{-1,0,1\} \right\},
		\]
		and assume, by way of contradiction, that $\mathbb{P}(A) > 0$. Therefore, for all $\omega \in A$,
		\[
		X(t,\omega) \sim \lambda(\omega) F^{-1}(t), \quad t \to \infty.
		\]
		Since $f$ is odd and obeys \eqref{eq.fasypres} and \eqref{eq.fasypresmu}, we have that 
		\[
		f(X(t,\omega)) \sim f(\lambda(\omega) F^{-1}(t)), \quad t \to \infty,
		\]
		in the case $\lambda\neq 0$ and, in the case where $\lambda=0$, we have $f(X(t,\omega))=o(f\circ F^{-1}(t)$ as $t\to\infty$.
		
		We consider the case where $\lambda=0$ first. Suppose this occurs on the event $A_0$, which we suppose to  be a zero probability event. In that case, since $t\mapsto f\circ F^{-1}(t)$ is in $L^1(0,\infty)$. Moreover, we have that $X(t)\to 0$ as $t\to\infty$, and by hypothesis we have $\sigma$ in $L^2(0,\infty)$ and hence $U(t)$ tend to a finite limit also. Therefore, proceeding as before, we get on $A_0$ 
		\[
		\Gamma(t) = X(t)-\int_t^\infty f(X(s))\,ds
		\] 
		Now, since $f(X(t))=o((f\circ F^{-1})(t))$ as $t\to\infty$ on $A_0$, 
		\[
		\int_t^\infty f(X(s))\,ds = o(F^{-1}(t)), \quad t\to\infty,
		\]
		and so on $A_0$, we have $\Gamma(t)=o(F^{-1}(t))$ as $t\to\infty$. 
		
		By Lemma~\ref{lemma.tailmartasy}, we have that
		\[
		\limsup_{t \to \infty} 
		\frac{|\Gamma(t)|}{\sqrt{2 \, \varphi\!\left(\int_t^\infty \sigma^2(s)\, ds\right)}} 
		= 1, 
		\quad \text{a.s.}
		\]
		Let the event on which this holds be $\Omega_6$. By hypothesis we have  
		\[
		\mu = \lim_{t \to \infty} \frac{\varphi\!\left(\int_t^\infty \sigma^2(s)\, ds \right)}{(F^{-1}(t))^2} > 0.
		\]
		Therefore, for $\omega \in A_0 \cap \Omega_6$ (which is an event of positive probability),
		\[
		\limsup_{t \to \infty} \frac{|\Gamma(t,\omega)|}{F^{-1}(t)}
		= \limsup_{t \to \infty} \frac{|\Gamma(t,\omega)|}
		{\sqrt{2 \varphi\!\left(\int_t^\infty \sigma^2(s)\, ds\right)}} 
		\cdot
		\frac{\sqrt{2 \varphi\!\left(\int_t^\infty \sigma^2(s)\, ds\right)}}{F^{-1}(t)}=\sqrt{2\mu},
		\]
		%	Since
		%	\[
		%	\lim_{t \to \infty} 
		%	\frac{\sqrt{2 \varphi\!\left(\int_t^\infty \sigma^2(s)\, ds\right)}}{F^{-1}(t)} 
		%	= \sqrt{2\mu},
		%	\]
		%	and
		%	\[
		%	\limsup_{t \to \infty} \frac{|\Gamma(t,\omega)|}
		%	{\sqrt{2 \varphi\!\left(\int_t^\infty \sigma^2(s)\, ds\right)}} = 1,
		%	\]
		%	we get
		%	\[
		%	\sqrt{2\mu} \cdot \limsup_{t \to \infty} 
		%	\frac{|\Gamma(t,\omega)|}{\sqrt{2 \varphi\!\left(\int_t^\infty \sigma^2(s)\, ds\right)}} 
		%	= \sqrt{2\mu} > 0.
		%	\]
		where we interpret the righthand side as $+\infty$ in the case $\mu=+\infty$. But this contradicts the conclusion that $X(t)=o(F^{-1}(t))$ on $A_0$, so we have that $\mathbb{P}[A_0]=0$. 
		
		We now turn to the cases where $\lambda\neq 0$. Consider $A_+=\{\omega:\lambda(\omega)>0\}$ and $A-=\{ \omega:\lambda(\omega)<0\}$. We have assumed $\mathbb{P}[A_+\cup A_-]>0$. Suppose that $\mathbb{P}[A_+]>0$; we will show that this leads to a contradiction. The case for $A_-$ follows similarly, exploiting the fact that $f$ is odd, so we only give a proof for $A_+$. 
		
		%We have already on $A_+$ that $f(X(t))\sim f(\lambda F^{-1}(t))=f(F^{-1}(t))$, and the latter function is in $L^1$. Therefore $t\mapsto f(X(t))$ is in $L^1(0,\infty)$ on $A_+$. provided 
		%$t\mapsto f(\lambda F^{-1}(t))$ is in $L^1(0,\infty)$. As in earlier calculations, we have 
		%\begin{equation} \label{eq.intrhs}
		%\int_0^t f(\lambda F^{-1}(s))\,ds = \int_{F^{-1}(t)}^1 \frac{f(\lambda u)}{f(u)}\,du\leq \int_{0}^1 \frac{f(\lambda u)}{f(u)}\,du.
		%\end{equation}
		%We have shown already that there exist constants $c>1$ and $a>0$ such that for all $\lambda>1$ we have 
		%\[
		%\limsup_{x\to\infty} \frac{f(\lambda x)}{f(x)} \leq c\lambda^a,
		%\]
		%and by monotonicity, for $\lambda\in (0,1]$, $f(\lambda x)
		%\leq f(x)$. Therefore, either way we get that the righthand side of 
		%\eqref{eq.intrhs} is finite, and therefore $t\mapsto f(X(t))$ is in $L^1(0,\infty)$. 
		Since $X(t)\to 0$ as $t\to\infty$ and $\sigma\in L^2(0,\infty)$, on $A_+$ we can once again write
		\[
		\Gamma(t) = X(t) - \int_t^\infty f(X(s)) \, ds,
		\]
		Since $\lambda=+1$, we see that both terms on the righthand side are asymptotic to $F^{-1}(t)$ as $t\to\infty$, so the righthand side is $o(F^{-1}(t))$ as $t\to\infty$. However, by reusing the argument above in the case $\lambda=0$, we have that 
		\[
		\limsup_{t\to\infty} \frac{|\Gamma(t)|}{F^{-1}(t)}=\sqrt{2\mu}>0,
		\] 
		giving a contradiction to the supposition that $\mathbb{P}[A_+]>0$, so $\mathbb{P}[A_+]=0$.
		As we noted before, a symmetric argument deals with $A_-$ and so $\mathbb{P}[A_-]=0$. Therefore, we have that $A=A_0\cup A_+\cup A_-$ is a zero probability event, completing the proof of part (ii).  
		
		For part (iii), if we assume there is an event $A'$ of positive probability on which $X(t)=O(F^{-1}(t))$ as $t\to\infty$, we can prove as before on $A'$ that $t\mapsto f(X(t))$ is in $L^1(0,\infty)$ and once again that on $A'$ that $\Gamma$ is given by 
		\[
		\Gamma(t) = X(t) - \int_t^\infty f(X(s)) \, ds.
		\]
		Then, using the argument from Theorem \ref{thm.bigOconverse} applied pathwise in $A'$, we see that on $A'$
		\[
		\int_t^\infty f(X(s)) \, ds = O(F^{-1}(t)), \quad t\to\infty.
		\]
		Since $X(t)=O(F^{-1}(t))$ as $t\to\infty$ on $A'$, we have that 
		$\Gamma(t)=O(F^{-1}(t))$ as $t\to\infty$ on $A'$. However, since 
		$\varphi(\int_t^\infty \sigma^2(s)\,ds)/F^{-1}(t)^2\to \infty$ as $t\to\infty$, we get a.s. that 
		\[
		\limsup_{t \to \infty} \frac{|\Gamma(t,\omega)|}{F^{-1}(t)}
		= \limsup_{t \to \infty} \frac{|\Gamma(t,\omega)|}
		{\sqrt{2 \varphi\!\left(\int_t^\infty \sigma^2(s)\, ds\right)}} 
		\cdot
		\frac{\sqrt{2 \varphi\!\left(\int_t^\infty \sigma^2(s)\, ds\right)}}{F^{-1}(t)}=+\infty,
		\]
		which contradicts what we deduced earlier i.e., that $\Gamma(t)=O(F^{-1}(t))$ as $t\to\infty$ on $A'$, where $A'$ is an event of positive probability. Therefore, the original supposition that $A'$ is an event of positive probability must be false, and so 
		$X(t)=O(F^{-1}(t))$ as $t\to\infty$ can only happen on an event whose probability is zero. From this (iii) follows. 
	\end{proof}

\end{document}
